% PPGEE - UFMG
%-------------------------------------------------------------------
% % Editando:
%     \documentclass[12pt,a4paper,titlepage]{book}
%     \usepackage[left=1.5cm,right=3.5cm,top=4cm,bottom=2.5cm]{geometry}
% Final:
    \documentclass[12pt,a4paper,titlepage,final]{book}
    \usepackage[left=2.5cm,right=2.5cm,top=4cm,bottom=2.5cm]{geometry}
% Rascunho:
    % \documentclass[12pt,a4paper,titlepage,draft]{book} \documentclass[12pt,a4paper,titlepage]{book}
    % \usepackage[left=1.5cm,right=3.5cm,top=4cm,bottom=2.5cm]{geometry}


% \documentclass[12pt,a4paper,titlepage,brazil,portugues]{book}
%--------------------------------------------------------------------
%--------------------------------------------------------------------
% Pacotes usados {{
\usepackage[utf8]{inputenc}
\usepackage[T1]{fontenc}
\usepackage{amsmath,amssymb,amsfonts,amstext,textcomp}         % permite simbolos matematicos
\usepackage{bm}                                                % negrito em letras gregas
\usepackage{mathrsfs}                                          % permite o uso de letras trabalhadas
\usepackage{amsthm}
\usepackage{pgf}
\usepackage{graphicx}
\usepackage{float} % Allows putting an [H] in \begin{figure} to specify the exact location of the figure
\usepackage{color}
\usepackage{xcolor}
\definecolor{myOrange}{rgb}{1,0.5,0}
% \definecolor{myOrange}{rgb}{1,0.5,0}
\usepackage{booktabs}                                          % Tabelas com qualidade de publicao
\usepackage{indentfirst}
\usepackage{rotating}
\usepackage{braket}
\usepackage[ruled,vlined]{algorithm2e}                         % Para algoritmos (melhor que o outro)
\usepackage{multirow}
\usepackage{fancyhdr}
\usepackage{makeidx}
\usepackage{setspace}
\usepackage[colorlinks,citecolor=blue]{hyperref}               % tirar [colorlinks] para imprimir
% \usepackage{autonum}                                         % Mostra numeracao somente em eq citadas
\usepackage{mathtools}
\mathtoolsset{showonlyrefs=true}                               % Mostra numeracao somente em eq citadas
\usepackage{ae}
\usepackage[below]{placeins}
\usepackage{flafter}
\usepackage{pxfonts}
% \usepackage{subfigure}
\usepackage{appendix}
\usepackage{listings}                                          % para cdigo fonte
\usepackage[symbol]{footmisc}                                  % Mudando marcadores do footnote

% Tikz
\usepackage{tikz}
\usepackage{pgfplots}
\usetikzlibrary{shapes,arrows,calc,automata,positioning}
\usetikzlibrary{plotmarks}    % marcas nos graficos
\usetikzlibrary{arrows.meta}  % estilos de setas

% \usetikzlibrary{external}
% \usepgfplotslibrary{external}
% % \tikzexternalize[prefix=tikz/]
% \tikzexternalize[optimize=false,prefix=Figs/]
% \tikzset{external/force remake=false}
\usepgfplotslibrary{patchplots}

\usepackage{soul}                                              % para highlighth (todonotes)
\usepackage{natbib}                                            % para usar citep, etc...
\usepackage{subfig}                                            % para subfiguras

\pgfplotsset{compat=newest}
%% the following commands are needed for some matlab2tikz features
\usepackage{grffile}   % Permite usar carc. esp. em nomes de arq
\usepackage{amsmath}   % Enhance-ments for improving mathematical formulas

\usepackage[
      textsize=scriptsize,
      colorinlistoftodos,
      % shadow,
      color=yellow,
      backgroundcolor=yellow,
      % obeyDraft
      obeyFinal
      ]{todonotes}
% \usepackage[
      % natbib=true,
      % % citestyle=authoryear,
      % citestyle=rusnat,
      % % bibstyle=authortitle,
      % % style=authoryear-comp,
      % hyperref=true,
      % backend=biber,
      % maxbibnames=99,
      % giveninits=true,
      % uniquename=init,
      % maxcitenames=2,
      % parentracker=true,
      % url=false,
      % doi=false,
      % isbn=false,
      % eprint=false,
      % backref=true]{biblatex}
% \usepackage[style    = abnt,           % Sistema alfabetico
          % % style    = abnt-numeric,   % Sistema numerico
          % % style    = abnt-ibid,      % Notas de referencia]{biblatex}
% ]{biblatex}
% \usepackage{mcode}
% \usepackage{amssymb}
% \usepackage{amsmath}
% \usepackage{caligra} %letras caligraficas
% \usepackage{caligr}%letras caligraficas
% \usepackage{dsfont}%barra escura para conjuntos
% \usepackage[dvips]{graphicx}
% \thispagestyle{empty}
% \usepackage{epsfig,latexsym}
% \usepackage{psfrag}
% \usepackage{epsfig}
% \usepackage[siunitx]{circuitikz} % fazer diagramas de blocos
% \usepackage{algorithm}
% \usepackage[noend]{algpseudocode}   % Para pseudo-codigos
% \usepackage[linesnumbered,ruled,vlined]{algorithm2e} % Para algoritmos (melhor que o outro)
% \usepackage[pdftex]{hyperref}
% \usepackage[hang,small]{caption}
% \usepackage{icomma}

% Para nao sobrar espaos em branco estranhos
% \widowpenalty=1000 \clubpenalty=1000
% \newcommand{\gepeto}[1]{\vspace*{12pt}\begin{singlespace}\hspace{2cm}\begin{small}\parbox{12cm}{#1}\end{small}\end{singlespace}\vspace*{20pt}}

\theoremstyle{definition}
\newtheorem{theorem}{Theorem}[section]
\newtheorem{corollary}{Corollary}[theorem]
\newtheorem{lemma}[theorem]{Lemma}
% \newtheorem{defi}{Definio}[section]
% \newtheorem{propo}{Proposio}[section]
% \newtheorem{coro}{Corolrio}[section]
% \newtheorem{exem}{Exemplo}[section]
% \newtheorem{lema}{Lema}[section]
% \newtheorem{consi}{Considerao}[section]


% Modificacoes em legenda
\setlength{\belowcaptionskip}{10pt}
% \addbibsource{referencias.bib}

% {{{ Modificas em legenda (quando ocorre quebra das palavras indesejadas)
% \hyphenation{na-tu-ra-is}
% \hyphenation{ge-ra-do}

% {{{ Definicao de comandos

\newcommand{\comment}[1]{}
\newcommand{\veq}{\vspace{0.2cm}}
\newcommand{\x}{\\[5pt]}
\newcommand{\C}[1]{\relax}

\newtheorem{teo}{Theorem}[section]
\newtheorem{defin}[teo]{Definition}
\newtheorem{nota}{Note}[section]

% }}}

% {{{ Novos operadores
\DeclareMathOperator*{\argmax}{argmax}
\DeclareMathOperator*{\argmin}{argmin}
% }}}

% {{{ Captulo personalizado
\makeatletter
\def\thickhrulefill{\leavevmode \leaders \hrule height 1ex \hfill \kern \z@}
\def\@makechapterhead#1{%
  \vspace*{10\p@}%
  {\parindent \z@ \centering \reset@font
        \thickhrulefill\quad
        \bfseries \@chapapp{} \thechapter
        \quad \thickhrulefill
        \par\nobreak
        \vspace*{10\p@}%
        \interlinepenalty\@M
        \hrule
        \vspace*{10\p@}%
        \Huge \bfseries #1\par\nobreak
        \par
        \vspace*{10\p@}%
        \hrule
    \vskip 70\p@    %100
  }}
\def\@makeschapterhead#1{%
  \vspace*{10\p@}%
  {\parindent \z@ \centering \reset@font
        \thickhrulefill
        \par\nobreak
        \vspace*{10\p@}%
        \interlinepenalty\@M
        \hrule
        \vspace*{10\p@}%
        \Huge \bfseries #1\par\nobreak
        \par
        \vspace*{10\p@}%
        \hrule
    \vskip 70\p@
 }
}

%--------------------------------------------------------------------
%Variacao da altura da linha
\renewcommand{\baselinestretch}{1.1}

%--------------------------------------------------------------------
%Coloca Rascunho e a data
\newcommand{\reviewtimetoday}[2]{\special{!userdict begin
/bop-hook{gsave 20 710 translate 45 rotate 0.8 0.8 0.8 0.8 0 setcmykcolor
/Times-Roman findfont 10 scalefont setfont 0 0   moveto (#1) show
0 -12 moveto (#2) show grestore}def end}}
% You can turn on or off this option.
\reviewtimetoday{\today}{Draft Version}

%--------------------------------------------------------------------
%Criar o index
\makeindex

%Devido ao fancyheading
\headheight 14.7pt

%--------------------------------------------------------------------
%Trocando os efeitos da virgula e ponto
% \DeclareMathSymbol{,}{\mathord}{letters}{"3B}
% \DeclareMathSymbol{.}{\mathpunct}{letters}{"3A}

% }}}
% -*- TeX-master: "Qualificacao.tex" -*-
%!TEX root = Qualificacao.tex

\newcommand{\norm}[1]{\left\lVert#1\right\rVert}

%%%%%%%%%%%%%%%%%%%%%%%%%%%%%%%%%%%%%%%%%%%%%%%%%%%%%%%%%%%%%%%%%%%%%%%
%                    Theorems, definitions, etc...                    %
%%%%%%%%%%%%%%%%%%%%%%%%%%%%%%%%%%%%%%%%%%%%%%%%%%%%%%%%%%%%%%%%%%%%%%%


\theoremstyle{plain}
\newtheorem{thm}{Theorem}[section]
\newtheorem{lem}[thm]{Lemma}
\newtheorem{prop}[thm]{Proposition}
\newtheorem*{cor}{Corollary}

\theoremstyle{definition}
\newtheorem{defn}{Definition}[chapter]
\newtheorem{assum}{Assumption}[chapter]
\newtheorem{conj}{Conjecture}[section]
\newtheorem{exmp}{Example}[section]

\theoremstyle{remark}
\newtheorem*{rem}{Remark}
\newtheorem*{note}{Note}

%%%%%%%%%%%%%%%%%%%%%%%%%%%%%%%%%%%%%%%%%%%%%%%%%%%%%%%%%%%%%%%%%%%%%%%
%                               Vectors                               %
%%%%%%%%%%%%%%%%%%%%%%%%%%%%%%%%%%%%%%%%%%%%%%%%%%%%%%%%%%%%%%%%%%%%%%%

\newcommand{\vet}[1]{\mathrm{\mathbf{#1}}}
\newcommand{\vx}{\vet{x}}
\newcommand{\vX}{\vet{X}}
% \newcommand{\vu}{\vet{u}}
\newcommand{\vu}{u}
\newcommand{\vy}{\vet{y}}
\newcommand{\vA}{\vet{A}}
\newcommand{\vB}{\vet{B}}
\newcommand{\vG}{\vet{G}}
\newcommand{\vH}{\vet{H}}
\newcommand{\vC}{\vet{C}}
\newcommand{\vD}{\vet{D}}
\newcommand{\vQ}{\vet{Q}}
\newcommand{\vR}{\vet{R}}
\newcommand{\vt}{\bm{\theta}}
\newcommand{\vf}{\vet{f}}
\newcommand{\vg}{\vet{g}}
\newcommand{\vK}{\vet{K}}
\newcommand{\vP}{\vet{P}}
\newcommand{\vv}{\vet{v}}
\newcommand{\vw}{\vet{w}}

\newcommand{\vtheta}{\bm{\theta}}
\newcommand{\vepsilon}{\bm{\epsilon}}
\newcommand{\vvarphi}{\bm{\varphi}}
\newcommand{\vPhi}{\bm{\Phi}}
\newcommand{\vPsi}{\bm{\Psi}}
\newcommand{\vxi}{\bm{\xi}}
\newcommand{\vmu}{\bm{\mu}}
\newcommand{\vpi}{\bm{\pi}}


%%%%%%%%%%%%%%%%%%%%%%%%%%%%%%%%%%%%%%%%%%%%%%%%%%%%%%%%%%%%%%%%%%%%%%%
%                        Virtual Signals                              %
%%%%%%%%%%%%%%%%%%%%%%%%%%%%%%%%%%%%%%%%%%%%%%%%%%%%%%%%%%%%%%%%%%%%%%%

\newcommand{\rv}{\bar{r}}
\newcommand{\ev}{\bar{e}}


%%%%%%%%%%%%%%%%%%%%%%%%%%%%%%%%%%%%%%%%%%%%%%%%%%%%%%%%%%%%%%%%%%%%%%%
%                            Other symbols                            %
%%%%%%%%%%%%%%%%%%%%%%%%%%%%%%%%%%%%%%%%%%%%%%%%%%%%%%%%%%%%%%%%%%%%%%%

\newcommand{\E}{\mathbb{E}}
\newcommand{\Prob}{\mathbb{P}}
\newcommand{\R}{\mathbb{R}}

\newcommand{\pt}{\tilde{\phi}}
\newcommand{\ph}{\hat{\phi}}

\newcommand{\Wh}[1]{\hat{W}_{#1}}
\newcommand{\Wht}[1]{\hat{W}_{#1}^T}
\newcommand{\Wt}[1]{\tilde{W}_{#1}}
\newcommand{\Wtt}[1]{\tilde{W}_{#1}^T}

\newcommand{\Wio}{W_i^*}
\newcommand{\Wiot}{W_i^{*T}}

\newcommand{\Wci}{\hat{W}_{ci}}
\newcommand{\Wai}{\hat{W}_{ai}}
\newcommand{\Wcit}{\hat{W}_{ci}^T}
\newcommand{\Wait}{\hat{W}_{ai}^T}

\newcommand{\laplace}[1]{\mathcal{L} \left\{ #1 \right\}}
\newcommand{\realiz}{$(A, B, C, D)$ }
\newcommand{\eqpar}[2]{\frac{\partial{#1}}{\partial{#2}}}

\newcommand{\xb}{\bar{x}}
% \newcommand{\norm}[1]{||#1||}
\newcommand{\Sb}{\bar{S}}

\newcommand{\Jio}{J_i^*}
\newcommand{\Jiz}{J_i^0}
\newcommand{\Jiho}{\hat{J}_i^*}

% \newcommand{\vpfa}{\mathbf{\phi}}

% Function approximations
\newcommand{\vpfa}{\bm{\varphi}}
\newcommand{\pfa}{\phi}
\newcommand{\Jfa}{\hat{J}}
\newcommand{\Jfaxp}{\hat{J}(\vx,\vpfa)}
\newcommand{\Jfaxpk}{\hat{J}(\vx_k,\vpfa_k)}
\newcommand{\Qfa}{\hat{Q}}
\newcommand{\Qfaxp}{\hat{Q}(\vx,u,\pfa)}

% Policy Gradient Methods
\newcommand{\apPG}{\nabla \hat{P}(\vt_k)}


\newcommand{\vuo}{\alpha_1^*}
\newcommand{\vuh}{\hat{\alpha}_1}
\newcommand{\vi}{\alpha_i}
\newcommand{\vio}{\alpha_i^*}
\newcommand{\viho}{\hat{\alpha}_i^*}
\newcommand{\vih}{\hat{\alpha}_i}
\newcommand{\vnh}{\hat{\alpha}_n}
\newcommand{\vnah}{\hat{\alpha}_{n-1}}
\newcommand{\vnahd}{\dot{\hat{\alpha}}_{n-1}}
\newcommand{\viah}{\hat{\alpha}_{i-1}}
\newcommand{\viahd}{\dot{\hat{\alpha}}_{i-1}}
\newcommand{\viaahd}{\dot{\hat{\alpha}}_{i-2}}


\newcommand{\tq}{\tilde{Q}}
\newcommand{\tqxu}{\tilde{Q}(x_k,u_k)}
\newcommand{\tqxuo}{\tilde{Q}(x_{k+1},u_{k+1})}



% Define a counter for the inserted todonotes.
\newcounter{todoListItems}
% \newcommand{\todoTrans}[2]{
% % Increment counter
% \addtocounter{todoListItems}{1}
% \todo[%
% caption={\protect\hypertarget{todo\thetodoListItems}{}[\thetodoListItems] #2}]
% {
% #1 \hfill
% \hyperlink{todo\thetodoListItems}{[\thetodoListItems]}
% }
% }

% highlight on todo notes
\makeatletter
   \if@todonotes@disabled
      \newcommand{\todohl}[2]{#1}
   \else
      \newcommand{\todohl}[2]{\texthl{#1}\todo{#2}}
   \fi
\makeatother
% \makeatletter
   % \if@todonotes@disabled
      % \newcommand{\todohl}[3]{#1}
   % \else
      % \newcommand{\todohl}[3]{\texthl{#1}\todo[#3]{#2}}
   % \fi
% \makeatother



% \graphicspath{{Figs/}}

\begin{document}


% {{{ DEFINICAO DE CORES
\newcommand{\blacksolidlinethin}{\raisebox{2pt}{\tikz{\draw[black,solid,line width=0.5pt](0,0) -- (8mm,0);}}}
\newcommand{\blacksolidline}{\raisebox{2pt}{\tikz{\draw[black,solid,line width=1pt](0,0) -- (8mm,0);}}}
\newcommand{\blackdottedline}{\raisebox{2pt}{\tikz{\draw[black,dotted,line width=1pt](0,0) -- (8mm,0);}}}
\newcommand{\blackdashdotline}{\raisebox{2pt}{\tikz{\draw[black,dashdotted,line width=1pt](0,0) -- (8mm,0);}}}
\newcommand{\bluesolidline}{\raisebox{2pt}{\tikz{\draw[blue,solid,line width=1pt](0,0) -- (8mm,0);}}}
\newcommand{\bluedashedline}{\raisebox{2pt}{\tikz{\draw[blue,dashed,line width=1pt](0,0) -- (8mm,0);}}}
\newcommand{\bluedashdotedline}{\raisebox{2pt}{\tikz{\draw[blue,dashdotted,line width=1pt](0,0) -- (8mm,0);}}}
\newcommand{\reddashedline}{\raisebox{2pt}{\tikz{\draw[red,dashed,line width=1pt](0,0) -- (8mm,0);}}}
\newcommand{\redsolidline}{\raisebox{2pt}{\tikz{\draw[red,solid,line width=1pt](0,0) -- (8mm,0);}}}
\newcommand{\reddottedline}{\raisebox{2pt}{\tikz{\draw[red,dotted,line width=1pt](0,0) -- (8mm,0);}}}
\newcommand{\greendashdottedline}{\raisebox{2pt}{\tikz{\draw[black!50!green,dashdotted,line width=1pt](0,0) -- (8mm,0);}}}
\newcommand{\magentadottedline}{\raisebox{2pt}{\tikz{\draw[magenta,dotted,line width=1pt](0,0) -- (8mm,0);}}}
% }}} FIM DE DEFINICAO DE CORES  ================================================

% % {{{ ELEMENTOS PRE-TEXTUAIS
%
% % %--------------------------------------------------------------------
% % %Pagina em algarismos romanos
% \pagenumbering{roman}
% % % \selectlanguage{brazil}
% % \selectlanguage{english}
% %
% %--------------------------------------------------------------------
% %Capa
% 
%CAPA
\begin{titlepage}

%Filiação
%--------------------------------------------------------------------
\begin{figure}[!ht]
\begin{minipage}[b]{0.7\linewidth}
\begin{tiny}
Laboratório de Modelagem, Análise e Controle de Sistemas Não-Lineares \\
Departamento de Engenharia Eletrônica\\
Universidade Federal de Minas Gerais\\
Av. Antônio Carlos 6627, 31270-901 Belo Horizonte, MG Brasil\\

Fone: +55 3499-4866 - Fax: +55 3499-4850 %\\
%emmendes@cpdeee.ufmg.br
\end{tiny}
\end{minipage}\hfill
\begin{minipage}[c]{0.3\linewidth}
\begin{flushright}
\vspace{-1.2cm}
\includegraphics[scale=1.37]{mac.png}
\end{flushright}
\end{minipage}
\end{figure}
%--------------------------------------------------------------------

%Título
%--------------------------------------------------------------------
\vspace{1.75cm}
\begin{center}
\thickhrulefill
\par\nobreak
\vspace*{10pt}%
\hrule
\vspace*{10pt}%
% {\bf \Huge \bfseries Controle não-linear baseado em dados: estrutura de controladores e uso de informação auxiliar}
% {\bf \Huge \bfseries Controle não linear baseado em dados: estrutura de controladores e uso de informação auxiliar}
{\bf \Huge \bfseries Non-linear Data-Driven Contol: Structure Selection of Controllers and Use of Auxiliary Information}
\vspace*{10pt}%
\hrule
\vspace*{40pt}%
%--------------------------------------------------------------------

%Autor
%--------------------------------------------------------------------
{\large \bf João Carlos Vilela de Castro}\\[0.5cm]
\end{center}
%--------------------------------------------------------------------


%Descrição
%--------------------------------------------------------------------
\vspace{1.75cm}
\begin{flushright}
\begin{minipage}{0.7\linewidth}
Texto de qualificação submetido à banca examinadora designada pelo Colegiado do
Programa de Pós-Graduação em Engenharia Elétrica da Universidade Federal
de Minas Gerais, como parte dos requisitos necessários à obtenção do grau
de Mestre em Engenharia Elétrica.\\
\end{minipage}
\end{flushright}
\vspace{1.6cm}



%--------------------------------------------------------------------

%Orientação
%--------------------------------------------------------------------
\begin{tabular}{ll}
  {\bf Orientador:} & Prof. Luis Antonio Aguirre, Ph.D.\\
    % &                        Prof., Dr.\\
\end{tabular}
\vspace{1.6cm}

%--------------------------------------------------------------------
%Cidade e Data
%--------------------------------------------------------------------
\centerline{Belo Horizonte, março de 2021}
%--------------------------------------------------------------------
\end{titlepage}

% \clearpage
% \thispagestyle{empty}
% \cleardoublepage
% %--------------------------------------------------------------------
% %
% % % %--------------------------------------------------------------------
% % % %Dedicatoria
% % % \input{dedicatoria.tex}
% % % \clearpage
% % % \thispagestyle{empty}
% % % \cleardoublepage
% % % %--------------------------------------------------------------------
% %
% %
% % % %--------------------------------------------------------------------
% % % %Agradecimentos
% % % \input{agradecimentos.tex}
% % % \clearpage
% % % \thispagestyle{empty}
% % % \cleardoublepage
% % % %--------------------------------------------------------------------
% %
% %
% % % %--------------------------------------------------------------------
% % % %Epigrafe
% % % \input{epigrafe.tex}
% % % \clearpage
% % % \thispagestyle{empty}
% % % \cleardoublepage
% % % %--------------------------------------------------------------------
% %
% %
% % % %%--------------------------------------------------------------------
% %Sumario
% \pagestyle{myheadings}
% \begin{spacing}{1.5}
% \tableofcontents
% \end{spacing}
% \clearpage
% \thispagestyle{empty}
% \cleardoublepage
% %
% %
% % % %%--------------------------------------------------------------------
% %Resumo
% \begin{spacing}{1}
% % -*- tex-master: "qualificacao.tex" -*-
%!tex root = qualificacao.tex

\chapter*{Resumo}
\addcontentsline{toc}{chapter}{Resumo}

\vspace{-2cm} 


Em procedimentos convencionais o projeto de controladores é feito a partir de um modelo matemático que representa o processo a se controlar. Outra abordagem, que difere em essência da convencional, com estudos crescentes nas últimas décadas, é a de projeto de controladores baseado em dados (DDC do inglês Data-Driven Control). No DDC, o projeto do controlador não faz uso direta ou indiretamente do modelo do processo e todo o projeto é feito a partir de dados amostrados diretamente do processo. Grande parte das técnicas DDC são métodos iterativos baseados principalmente no método do gradiente para minimizar algum índice de custo. Contudo algumas técnicas, em especial a VRFT, do inglês Virtual Reference Feedback Tuning, permitem, por um procedimento em batelada, realizar a minimização deste índice a partir de técnicas usuais no âmbito da identificação de sistemas. Em geral estes procedimentos são feitos fixando-se uma estrutura para o controlador, e a partir de métodos de identificação e otimização procura-se pelo melhor conjunto de parâmetros para o controlador que resulte em um comportamento próximo ao definido por um modelo de referência. Nesta pesquisa procura-se por uma abordagem alternativa, em que a melhor estrutura do controlador seja selecinada no processo de sintonia do controlador.  Neste sentido tem-se buscado um método para auxilio na seleção da melhor estrutura do controlador a partir de uma estratégia de controle VRFT para sistemas não lineares. Esta seleção é feita por uma abordagem aleatorizada de seleção de estruturas de modelos recente no âmbito de identificação de sistemas e que é, neste trabalho, adaptada para lidar com identificação de controladores. Como um texto de qualificação, alguns resultados prévios alcançados são estudados, e propostas para continuidade do trabalho são apresentadas. 

\textbf{Palavras-chave:} Controle baseado em dados; Controle livre de modelo; Aprendizado por reforço; Seleção de estruturas; Sistemas não lineares.

% \clearpage
% \thispagestyle{empty}
% \cleardoublepage
% % % %%--------------------------------------------------------------------
% %
% %
% % % %%--------------------------------------------------------------------
% %Abstract
% % -*- tex-master: "../Qualificacao.tex" -*-
%!tex root = qualificacao.tex

\chapter*{Abstract}
\addcontentsline{toc}{chapter}{Abstract}

\vspace{-2cm}

In conventional procedures, the design of controllers is based on a mathematical model that represents the process to be controlled. Another approach, which differs in essence from the conventional one, with increasing studies in the last decades, is the design of data-driven controllers (DDC). In the DDC approach, the controller project does not make direct or indirect use of the process model and the entire project is made from data sampled directly from the process. Most of the DDC techniques are iterative methods based mainly on the gradient method to minimize some cost index. However, some techniques, especially the Virtual Reference Feedback Tunning (VRFT), allows, through a batch procedure, to minimize this index using the usual techniques in the scope of systems identification. In general, these procedures are done by setting a structure for the controller, and based on identification methods, the best set of parameters for the controller is identified, aiming to results in a behavior close to that defined by a reference model. In this research, an alternative approach is proposed, in which the objective is to select the best structure of the controller, before the process of tuning or identification of the controller's parameters. In this sense, a method has been developed to assist in the selection of the best controller structure based on a VRFT control strategy for non-linear systems. This selection is made with a modified version of a recent published randomized  model structure selection approach used in model identification. An analise of the structure selection problem is made in sense of controller structure selection and parameter identification. As a qualification text, some preliminary results achieved are studied, and proposals for further work are presented.


proposta

\textbf{Keywords:} Data-Driven Control; Model-Free Control; Reinforcement Learning; Structure Selection; Nonlinear Systems.

% Em procedimentos convencionais o projeto de controladores é feito a partir de um modelo matemático que representa o processo a se controlar.
% Uma outra abordagem, que difere em essência da convencional, com estudos crescentes nas últimas décadas, é a de projeto de controladores baseado em dados (DDC do inglês Data-Driven Control).
% No DDC, o projeto do controlador não faz uso direta ou indiretamente do modelo do processo e todo o projeto é feito a partir de dados amostrados diretamente do processo.
% Grande parte das técnicas DDC são métodos iterativos baseados principalmente no método do gradiente para minimizar algum índice de custo.
% Contudo algumas técnicas, em especial a técnica VRFT (do inglês Virtual Reference Feedback Tuning), permitem, por um procedimento em batelada, realizar a minimização deste índice a partir de técnicas usuais no âmbito da identificação de sistemas.
% No contexto de identificação de sistemas é consenso que a estrutura do modelo – quais regressores o compõem – tem forte influência no desempenho dinâmico.
% Neste sentido esta pesquisa de doutorado tem buscado um método para auxilio na seleção da melhor estrutura do controlador a partir de uma estratégia de controle VRFT para sistemas não-lieares.
% Esta seleção é feita por uma abordagem aleatorizada de seleção de estruturas de modelos já conhecida no âmbito de identificação de sistemas, mas que é, neste trabalho, adaptada para lidar com identificação de controladores.
% Afim de reduzir custo computacional e tornar o procedimento mais viável, o uso de técnicas de aprendizado por reforço é considerado.
% Por fim, o uso de informações auxiliares, que tem mostrado benefícios no contexto de identificação de sistemas, é analisado no contexto da identificação de controladores, por meio de restrições que levam em conta características previamente conhecidas no projeto do controlador.




% \clearpage
% \thispagestyle{empty}
% \cleardoublepage
% \end{spacing}
% % % %%--------------------------------------------------------------------
% %
% % %%--------------------------------------------------------------------
% % %Lista de figuras
% \begin{sloppypar}
% \listoffigures
% \end{sloppypar}
% \addcontentsline{toc}{chapter}{Lista of Figures}
% \clearpage
% \thispagestyle{empty}
% \cleardoublepage
% % %%--------------------------------------------------------------------
% %
% % %%Lista de Tabelas
% \begin{spacing}{1}
% % \listoftables
% % \addcontentsline{toc}{chapter}{List of Tables}
% % \clearpage
% % \thispagestyle{empty}
% % \cleardoublepage
% % %
% %%--------------------------------------------------------------------
% %%Simbolos
%   % -*- TeX-master: "Qualificacao.tex" -*-
%!TEX root = Qualificacao.tex

\chapter*{List of Symbols}
\addcontentsline{toc}{chapter}{List of Symbols}
\vspace{-2.0cm}
\section*{Chapter 1}
      \footnotesize
      \begin{tabular}{lp{10cm}}
      $y_{k}$                  & Output signal $ \in \mathbb{R}^p$, at time $k$;
      % $\mathbb{R}^{l}$;\\
      % $x_{k}$                 & Vetor de estados no instante $k$ $\in$
      % $\mathbb{R}^{n}$;\\
      % $A$                     & Matriz dinâmica do sistema $\in$
      % $\mathbb{R}^{n\times n}$ ;\\
      % $B$                     & Matriz de entrada $\in$ $\mathbb{R}^{n\times
      % m}$;\\
      % $C$                     & Matriz de saída $\in$ $\mathbb{R}^{l\times n}$;\\
      % $D$                     & Matriz de transmissão direta $\in$ $\mathbb{R}^{l\times m}$;\\
      % $w_{k}$                 & Vetor de ruído de processo $\in$
      % $\mathbb{R}^{n}$;\\
      % $v_{k}$                 & Vetor de ruído de medição  $\in$ $\mathbb{R}^{l}$ ;\\
      % $Q$                     & Matriz de covariância da sequência de ruído $w_{k}$, \textbf{E}($w_{k}w_{k}^T$) $\in$ $\mathbb{R}^{n\times n}$;\\
      % $R$                     & Matriz de covariância da sequência de ruído $v_{k}$, \textbf{E}($v_{k}v_{k}^T$) $\in$ $\mathbb{R}^{l\times l}$;\\
      % $S$                     & Matriz de covariância das sequências de ruído $w_{k}$ e $v_{k}$, \textbf{E}($w_{k}v_{k}^T$) $\in$ $\mathbb{R}^{n\times l}$;\\
      % $\mathbb{R}$            & Números reais;\\
      % $\mathbb{R}^{n}$        &$\mathbb{R}^{n \times 1}$ (vetores reais tipo coluna);\\
      % $\in$                   & Pertence (é um elemento de);\\
      % \textbf{E[$\bullet$]}   & Esperança matemática;\\
      % $(\bullet)^{T}$         & Transposição de vetores ou matrizes;\\
      % $N$                     & Número  de medições;\\
      % $n$                     & Ordem do modelo;\\
      % $\hat{x}$                & Valor estimado de $x$;\\
      % $||A||_F$               & Norma de Frobenius da matriz $A$, $\sqrt{tr\;A^{*}A}$;\\
      % $tr\;A$                 & Traço de $A$;\\
      % $A^{*}$                 & $\overline{A}^{T}$  Conjugado transposto de $A$; \\
      % $\mathbb{N}$            & Espaço dos números naturais; \\
      % $U_{0|2i-1}$            & Matriz em blocos de Hankel da entrada $\in$ $\mathbb{R}^{2mi\times j}$;\\
      % $i$ e $j$               & Índices definidos pelo usuário $\in$
      % $\mathbb{N}$, tal que, $i$ $\geq$ $n$ e $j$ $\gg$ $i$;\\
      % $U_{p}$ e $U_{f}$       & Notação utilizada na matriz em blocos de Hankel da entrada para dividí-la\\
                              % & em dados passados ($p$) e futuros ($f$);\\
      % $u^{\bullet}$           & Faz referência a uma dada entrada;\\
      % $Y_{0|2i-1}$            & Matriz em blocos de Hankel da saída;\\
      % $Y_{p}$ e $Y_{f}$       & Notação utilizada na matriz em blocos de Hankel da saída para dividí-la \\
                              % & em dados passados ($p$) e futuros ($f$);\\
      % $(\bullet)^{+}$         & Adicionar um bloco linha em $(\bullet)$;\\
      % $(\bullet)^{-}$         & Retirar um bloco linha em $(\bullet)$;\\
      % $\square$               & Fim de demonstração; \\
      % $W_{p}$                 & Matriz em blocos de Hankel definida pelo empilhamento de entradas e saída;\\
      % $X_{i}$                 & Sequências de estados $\in \mathbb{R}^{n\times j}$;\\
      % $X_{p}$                 & Sequências de estados passados;\\
      % $X_{f}$                 & Sequências de estados futuros;\\
      \end{tabular}

\section*{Chapter 2}
      \footnotesize
      \begin{tabular}{lp{11cm}}
            $N$            & number of sampled data;\\
            $N_{p}$            & number of sampled models;\\
            % $\vy$          & vector of output data $\in$ $\mathbb{R}^{N}$;\\
            $\mathcal{I}$ & \\
            $\mathcal{I}_{j}^+$ & \\
            $\mathcal{I}_{j}^-$ & \\
             $\mathcal{J}$ & performance index \\
             $\mathscr{M}$ & universe set of all considered possible models \\
             $\tilde{\mathscr{M}}$ & \\
             $\mathscr{R}$ & set of all model candidate regressors \\
             $f^*$ & \\
             $\tilde{f}$ & \\
             $\E[\cdot]$ & expected value operator \\
      \end{tabular}

\section*{Chapter 3}
      \footnotesize
      \begin{tabular}{p{4cm}p{10.5cm}}
      \end{tabular}

\section*{Chapter 4}
      \footnotesize
      \begin{tabular}{lp{10cm}}
            $u_{k}$                  & vector input signal $\in$ $\mathbb{R}^{m}$, at time $k$; \\
            $\vu_{k}$                  & vector input signal $\in$ $\mathbb{R}^{m}$, at time $k$; \\
      \end{tabular}

\section*{Chapter 5}
      \footnotesize
      \begin{tabular}{lp{11cm}}
      \end{tabular}


%   \clearpage
%   \thispagestyle{empty}
%   \cleardoublepage
% %%--------------------------------------------------------------------
% %
% % %%%Abreviaturas
% %%--------------------------------------------------------------------
%   \input{abreviations.tex}
%   \clearpage \thispagestyle{empty}
%   \cleardoublepage
% \end{spacing}
% % %--------------------------------------------------------------------
% %
% % %Comecando os capitulos
% % %--------------------------------------------------------------------
% \clearpage
% \pagestyle{myheadings}
% \pagenumbering{arabic}
% % %--------------------------------------------------------------------
%
% % }}} FIM DE ELEMENTOS PRE-TEXTUAIS  ================================================

% {{{ Defincao de cabecalhos -------------------------------------------
\pagestyle{fancy}
\renewcommand{\chaptermark}[1]{\markboth{\thechapter\ #1}{}}
\renewcommand{\sectionmark}[1]{\markright{\thesection\ #1}{}}
\fancyhf{}
\fancyhead[LE,RO]{\thepage}
\fancyhead[LO]{\rightmark}
\fancyhead[RE]{\leftmark}
\renewcommand{\headrulewidth}{0.5pt}
\renewcommand{\footrulewidth}{0.0pt}
\addtolength{\headheight}{2.5pt}                % 2.5pt
\fancypagestyle{plain}{
     \fancyhead{}
     \renewcommand{\headrulewidth}{0pt}
     }
% }}} ------------------------------------------------------------------




\begin{algorithm}[htpb]
  \caption{RaMSS algorithm}\label{alg:RaMSS}
  $\bm{y},N_p,m,\mu_{\min},\mu_{\max},K,\mathscr{R}=\{\phi_{1},\dots,\phi_{m}\}$ \\
  \While{$iter < iter_{\max}$} 
   {
     $ \bm{\mu} \gets \bm{\mu}_0$ \\
     \For{$i=1:N_p$}
      {
         $\tau \gets 0$; $\bm{\psi}(k) \gets [\ ]$ \\ 
         \For{$j=1:m$}
         {
           $r_j \gets \text{Be}(\mu_j)$ \\
            \If{$r_j = 1$}
            {
            $\bm{\psi}(k) \gets \begin{bmatrix} \bm{\psi}^T(k) & \phi_j \end{bmatrix}^T $ \\
            $\tau \gets \tau + 1$ \\
            }
         }
         $\hat{\bm{\theta}}$ \gets LS \\
         \For{$h=1:\tau$}
         {
           $\tilde{\bm{\psi}}(k) \gets \text{non-redundant}(\bm{\psi(k)})$ \\
         }
         $\hat{\bm{y}} \gets \text{Predict}(\tilde{\bm{\psi}}(k))$ \\
         $\mathcal{J}_i \gets e^{-K\cdot MSPE(\bm{y},\bm{\hat{y}})} $
      }
      \For{$j=1:m$}
      {
         $\mathcal{J}^{+} \gets 0$; $ \mathcal{J}^{-} \gets 0$; $n^{+} \gets 0$; $n^{-} \gets 0$; \\
         \For{$i=1:N_p$}
         {
            \If{$\phi_j(k) \in \tilde{\bm{\psi}}(k)$}
            {
               $\mathcal{J}^{+} \gets \mathcal{J}^{+} + \mathcal{J}_i$ \\
               $n^{+} \gets n^{+} + 1$
            } \Else
            {
               $\mathcal{J}^{-} \gets \mathcal{J}^{-} + \mathcal{J}_i$ \\
               $n^{-} \gets n^{-} + 1$
            }
          }
            $\mathcal{I}_j \gets \left( \frac{\mathcal{J}^{+}}{\max(n^+,1)} - \frac{\mathcal{J}^{-}}{\max(n^-,1)} \right) $ \\
         $\mu_j \gets \mu_j + \gamma \mathcal{I}_j$ \\
         $\mu_j \gets \max \left( \min(\mu_j, \mu_{\max}), \mu_{\min} \right) $ \\
         
      }
   }
\end{algorithm}%[H]


%--------------------------------------------------------------------
\end{document}

% vim: set sw=4 ts=4 sts=2 tw=80 ft=tex 

